\subsection{Setup}
\D{series, absolute convergence}{
$\sum_{n=1}^\infty a_n:=\lim_{N\to\infty}s_N$, $s_N=\sum_{n=1}^N a_n$ (a series IS the sequence of partial sums).
\hd{absolutely convergent:} $\sum|a_n|<\infty$. \textbf{abs.\ conv.\ \ra\ conv.} (not conversely: $\sum\frac{(-1)^n}n$).
\hd{conditionally convergent:} converges but not absolutely.
}

\subsection{Decision tree: does $\sum a_n$ converge?}
\R{run these in order, stop at the first hit}{
\begin{rcp}
\item \hd{$a_n\not\to0$?} \ra\ \textbf{diverges}, done \de{Trivialkriterium}. (Always check this first, it is free.)
\item \hd{Known shape?} geometric $\sum q^n$ conv.\ $\iff|q|<1$ (value $\frac{1}{1-q}$ from $n=0$);
$p$-series $\sum\frac1{n^p}$ conv.\ $\iff p>1$; $\sum\frac1{n\ln^p n}$ conv.\ $\iff p>1$.
\item \hd{Factorials or $c^n$?} \ra\ \textbf{ratio test} $\lim\left|\frac{a_{n+1}}{a_n}\right|=:q$. $q<1$ conv.\ (abs.), $q>1$ div., $q=1$ useless \ra\ next test.
\item \hd{$n$-th powers ($a_n=b_n^{\,n}$, or $n$ in the exponent)?} \ra\ \textbf{root test} $\lim\sqrt[n]{|a_n|}=:q$, same rule. Root test is \emph{stronger} than ratio.
\item \hd{Looks like a known series?} \ra\ \textbf{comparison} \de{Majorante/Minorante}: $0\le a_n\le b_n$, $\sum b_n$ conv.\ \ra\ $\sum a_n$ conv.; $a_n\ge b_n\ge0$, $\sum b_n$ div.\ \ra\ $\sum a_n$ div. Compare only \emph{leading powers}: $\frac{3n^2+n}{n^4-1}\sim\frac3{n^2}$ \ra\ conv.
\item \hd{Alternating $(-1)^n b_n$?} \ra\ \textbf{Leibniz}: $b_n\ge0$, \emph{monotonically} decreasing, $b_n\to0$ \ra\ converges. (Remainder $|s-s_N|\le b_{N+1}$.)
\item \hd{$a_n=f(n)$, $f\ge0$ decreasing?} \ra\ \textbf{integral test}: $\sum f(n)$ and $\int_1^\infty f\dx$ share the same behaviour.
\item \hd{Telescoping} $a_n=b_n-b_{n+1}$ \ra\ $s_N=b_1-b_{N+1}$; split by partial fractions first.
\end{rcp}
}
\X{test traps}{
\begin{bul}
\item Ratio/root $q=1$ gives \textbf{no information} ($\sum\frac1n$ and $\sum\frac1{n^2}$ both have $q=1$).
\item Leibniz needs \emph{monotone} $b_n$ --- $b_n=\frac{2+(-1)^n}{n}\to0$ but not monotone.
\item Comparison needs $a_n\ge0$; otherwise apply it to $|a_n|$ (\ra\ absolute convergence).
\item Convergence is unaffected by finitely many terms / by the starting index.
\item $\sum a_n$ conv.\ and $\sum b_n$ conv.\ \ra\ $\sum(a_n+b_n)$ conv. But conv.\ $\times$ conv.\ termwise says nothing.
\end{bul}
}
\E{run the tree on five series}{
\begin{bul}
\item $\sum\frac{n^2}{2^n}$: $c^n$ \ra\ ratio $=\frac{(n+1)^2}{2n^2}\to\frac12<1$ \ra\ \textbf{conv.}
\item $\sum\frac{n!}{n^n}$: ratio $=\frac{(n+1)!n^n}{n!(n+1)^{n+1}}=\left(\frac{n}{n+1}\right)^n\to\frac1e<1$ \ra\ \textbf{conv.}
\item $\sum\left(\frac{n}{2n+1}\right)^n$: $n$-th power \ra\ root $=\frac n{2n+1}\to\frac12<1$ \ra\ \textbf{conv.}
\item $\sum\frac{2n+3}{n^2+7}$: leading behaviour $\frac2n$ \ra\ compare with $\sum\frac1n$ \ra\ \textbf{div.}
\item $\sum\frac{(-1)^n}{\sqrt n}$: $\frac1{\sqrt n}\downarrow0$ \ra\ Leibniz \ra\ conv.; but $\sum\frac1{\sqrt n}$ div.\ \ra\ \textbf{conditionally conv.}
\end{bul}
}
\F{rearrangement and products}{
\begin{bul}
\item \hd{Riemann rearrangement:} conditionally convergent \ra\ reordering can produce \emph{any} value in $\Rl\cup\{\pm\infty\}$.
\item \hd{Dirichlet:} absolutely convergent \ra\ every reordering gives the same sum.
\item \hd{Cauchy product} $c_n=\sum_{k=0}^n a_kb_{n-k}$: if at least one of the two converges absolutely then $\sum c_n=(\sum a_n)(\sum b_n)$ (Mertens).
\end{bul}
}
\F{values to know}{
$\sum_{n\ge0}q^n=\frac1{1-q}$, $\sum_{n\ge1}nq^{n-1}=\frac1{(1-q)^2}$ $(|q|<1)$;
$\sum_{n\ge0}\frac{x^n}{n!}=e^x$; $\sum_{n\ge1}\frac{(-1)^{n+1}}{n}=\ln2$;
$\sum_{n\ge1}\frac1{n^2}=\frac{\pi^2}6$; $\sum_{n\ge1}\frac1{n(n+1)}=1$.
}

\subsection{Power series $\sum a_n(x-x_0)^n$}
\R{"for which $x$ does it converge?"}{
\begin{rcp}
\item Read off the centre $x_0$ (in $(x-2)^n$: $x_0=2$) and the coefficient $a_n$ (everything without $(x-x_0)^n$).
\item \hd{Radius:} $\dfrac1R=\lim\left|\dfrac{a_{n+1}}{a_n}\right|$ (factorials!) \quad or \quad $\dfrac1R=\limsup\sqrt[n]{|a_n|}$ (powers). $R=\infty$ if the limit is $0$; $R=0$ if it is $\infty$.
\item Converges absolutely on $|x-x_0|<R$, diverges on $|x-x_0|>R$.
\item \hd{Endpoints $x=x_0\pm R$ separately!} Plug them in and use \S3.2 (usually Leibniz or $p$-series).
\item Answer as an interval, e.g. $[x_0-R,\ x_0+R)$.
\end{rcp}
}
\E{$\sum_{n\ge1}\frac{n^n(x-2)^n}{(2n)!}$}{
$a_n=\frac{n^n}{(2n)!}$ (factorial \ra\ ratio test).
$\left|\frac{a_{n+1}}{a_n}\right|=\frac{(n+1)^{n+1}}{(2n+2)!}\cdot\frac{(2n)!}{n^n}
=\frac{(n+1)\left(1+\frac1n\right)^n}{(2n+1)(2n+2)}\xrightarrow[n\to\infty]{}\frac{e\,(n+1)}{4n^2}\to0.$
So $1/R=0$, $R=\infty$: \textbf{converges for all $x\in\Rl$.}
}
\F{what power series give you for free}{
Inside $|x-x_0|<R$: the sum is $C^\infty$, and you may \textbf{differentiate and integrate termwise}; $R$ does not change.
Convergence is \emph{uniform} on every closed $[x_0-r,x_0+r]$ with $r<R$ (not necessarily on the whole open interval).
Coefficients are unique: $a_n=\frac{f^{(n)}(x_0)}{n!}$.
}
